\documentclass[11pt]{amsart}

\usepackage[T1]{fontenc}
\usepackage{lmodern}
\usepackage{microtype}
\usepackage{mathtools,amssymb,amsthm}
\usepackage[hidelinks]{hyperref}

\hypersetup{
  pdftitle={The reverse degree-polynomial range clause for the 2 x n Miura-ori
            flip graph fails for every a >= 7}
}

\newtheorem{theorem}{Theorem}
\newtheorem{proposition}[theorem]{Proposition}
\newtheorem{lemma}[theorem]{Lemma}
\newtheorem{corollary}[theorem]{Corollary}
\theoremstyle{remark}
\newtheorem*{remark}{Remark}

\newcommand{\OFG}{\mathrm{OFG}}
\newcommand{\N}{\mathbb{N}}

\title[The reverse degree-polynomial range clause fails for every $a\ge7$]
{The Reverse Degree-Polynomial Range Clause for the
 $2\times n$ Miura-ori Flip Graph Fails for Every $a\ge 7$}

\date{}

\subjclass[2020]{05C07, 05A15, 52C99}
\keywords{origami flip graph, Miura-ori, degree sequence, quasi-polynomiality,
generating function, counterexample}

\begin{document}

\begin{abstract}
Let $v_n^d$ be the number of vertices of degree $d$ in the origami flip graph
$\OFG(M_{2,n})$ of the $2\times n$ Miura-ori, and put $f_a(n):=v_n^{2n-a}$.
A labelled theorem of Christensen, Hull, O'Neil, Pappano, Ter-Saakov and Yang,
printed with no proof, asserts that $f_a$ agrees with a polynomial of degree
$\lfloor a/2\rfloor$ on the whole of $n\ge\lceil a/2\rceil+1$, for every
$a\in\N_0$. We refute the range clause. The witness at $a=7$ consists of five
integers printed in that paper's own degree-distribution table,
$f_7(5),\dots,f_7(9)=12,88,296,680,1288$, whose fourth forward difference is
$4\ne 0$ although the window lies inside the asserted range. A
generating-function argument then shows the failure is an infinite family:
writing $k_a=\lfloor a/2\rfloor$ and
$G_a=(1-x)^{k_a+1}\sum_{n\ge1}f_a(n)x^n$, we prove $G_a$ is a polynomial of
degree $3\lfloor a/2\rfloor-1$ or $3\lfloor a/2\rfloor$ according to the parity
of $a$, with leading coefficient exactly $4(-1)^{\lfloor a/2\rfloor+1}$ and
$G_a(1)>0$. Hence the asserted statement is false for every $a\ge7$, with its
last obstruction at $n=a-2$. Granting the source's own exclusion of the
degree-$2$ cell, $a=0$ and $2\le a\le6$ do satisfy the asserted range; we assert
nothing at the degenerate index $a=1$. What survives
is proved rather than refuted: polynomiality, and degree exactly
$\lfloor a/2\rfloor$, hold for all $a\ge2$, and the repaired hypothesis is
$n\ge\max(\lfloor a/2\rfloor+2,\,a-1)$.
\end{abstract}

\maketitle

\section{The statement, with its locator}

Let $M_{2,n}$ be the $2\times n$ Miura-ori crease pattern and $\OFG(M_{2,n})$
its origami flip graph, whose vertices are the flat-foldable mountain--valley
assignments of $M_{2,n}$ and whose edges join two assignments differing by a
face flip. Following Christensen, Hull, O'Neil, Pappano, Ter-Saakov and
Yang~\cite{CHOPTY}, $v_n^d$ denotes the number of vertices of degree $d$ in
$\OFG(M_{2,n})$, and we abbreviate the \emph{reverse diagonal}
\[
  f_a(n) \ :=\ v_n^{\,2n-a}\qquad (a\in\N_0).
\]

Our target is the theorem at lines 1256--1258 of the arXiv e-print \LaTeX{}
source \texttt{Miuradraft-v2.tex} of~\cite{CHOPTY}; we call it the
\emph{reverse-diagonal theorem}
below. We quote from the source file rather than the PDF so that the locator is
exact.

\begin{quote}\footnotesize
\begin{verbatim}
\begin{theorem} \label{thm:reversepolynomials}
The explicit formula for $v_n^{2n-a}$ where $n \ge
\left\lceil \frac{a}{2} \right\rceil+1$ is a polynomial
of $n$ with degree $\left\lfloor \frac{a}{2}
\right\rfloor$ for all $a \in \N_0$.
\end{theorem}
\end{verbatim}
\end{quote}

\noindent
(The source lines have been re-wrapped to fit this page; not a character of
the statement is altered.) The sentence that introduces the theorem, at line
1254, reads:

\begin{quote}
``The data in Table \texttt{\textbackslash ref\{tab:reversepolynomialdata\}}
suggests the following Theorem, which can be proven using
Proposition~\texttt{\textbackslash ref\{prop:recurrences\}} and induction.''
\end{quote}

\noindent
The environment then closes with no proof. Two further lines of the same file
are used below and are quoted here because our argument rests on them.

\begin{itemize}
\item Line 1236 fixes the index set and states a carve-out in one paragraph:
the reverse sequences are collected ``for all $a \in \N_0$, where
$\N_0=\N\cup\{0\}$'', ``this time excluding the edge case where there are
exactly $4$ vertices labeled with $2n-a$ (i.e.\ $2n-a=2$)''.
\item Line 1204 records the anomaly $v_1^2=2$, against $v_m^2=4$ for $m\ge2$.
\item The source's recurrence proposition of lines 1189--1202, which it
\emph{proves} there, gives
\[
 b_n^d=v_{n-1}^{d-2},\qquad
 w_n^d=v_{n-1}^{d}+w_{n-1}^{d-1}+v_{n-2}^{d-2},
\]
\[
 v_n^d=v_{n-1}^{d}+w_{n-1}^{d-1}+v_{n-1}^{d-2}+v_{n-2}^{d-2}.
\]
\end{itemize}

\noindent
Eliminating $w$ from the last two identities (as~\cite{CHOPTY} itself does at
line 1219) puts the recurrence in pure-$v$ form:
\begin{equation}\label{eq:P}
  v_n^d \;=\; v_{n-1}^{d}+v_{n-1}^{d-1}+v_{n-1}^{d-2}
            +v_{n-2}^{d-2}-v_{n-2}^{d-3}
  \qquad (n\ge 3),
\end{equation}
with the conventions of~\cite{CHOPTY}: $v_m^e=0$ for $e<2$ or $e>2m$,
$v_1^2=2$, and $v_m^2=4$ for $m\ge2$. We take
\eqref{eq:P} as given; it is a consequence of a proposition the source proves.

\medskip\noindent\textbf{Two readings of the admissible start.}
Since $2n-a=2$ has an integer solution only for even $a$, namely at
$n=a/2+1=\lceil a/2\rceil+1$, the theorem's literal first admissible $n$ is
itself the cell excluded at line 1236 when $a$ is even. Read literally the
theorem therefore already fails at $a=6$, where the claimed degree is
$\lfloor6/2\rfloor=3$ and the first admissible index is $n=4$: taking
$f_6(4)=v_4^2=4$ from the convention at line 1204,
\[
 f_6(4),\dots,f_6(8)\;=\;4,\;36,\;128,\;292,\;544,
\]
whose fourth forward difference is $544-4\cdot292+6\cdot128-4\cdot36+4=4\ne0$.
That is not what we rest on. Throughout what follows we
instead grant the source its own exclusion clause and work from
\begin{equation}\label{eq:s}
  s(a)\;:=\;\Big\lfloor \tfrac a2\Big\rfloor+2,
\end{equation}
which for \emph{odd} $a$ equals the literal bound $\lceil a/2\rceil+1$ and for
\emph{even} $a$ is one larger. That makes the refutation harder, and it makes
the $a=7$ witness below independent of any quantifier quibble. Consistently
with \eqref{eq:s}, the printed rows of the source's reverse-polynomial data
table begin at
$n=\lfloor a/2\rfloor+2$ in all six of the rows $a=0,\dots,5$ it displays.

\section{The witness at \texorpdfstring{$a=7$}{a=7}}

At $a=7$ the claimed degree is $\lfloor7/2\rfloor=3$, the claimed range is
$n\ge\lceil7/2\rceil+1=5=s(7)$, and the exclusion clause is vacuous, since
$2n-7=2$ has no integer solution.

\begin{theorem}\label{thm:a7}
The reverse-diagonal theorem of~\cite{CHOPTY} quoted above is false at $a=7$.
\end{theorem}

\begin{proof}
The five values
\begin{equation}\label{eq:witness}
  f_7(5),f_7(6),f_7(7),f_7(8),f_7(9)\;=\;12,\;88,\;296,\;680,\;1288
\end{equation}
are cells of the source's own degree-distribution table. Their Newton
forward-difference table is
\[
\begin{array}{lccccc}
 f      & 12 & 88 & 296 & 680 & 1288\\
 \Delta^1&    & 76 & 208 & 384 & 608\\
 \Delta^2&    &    & 132 & 176 & 224\\
 \Delta^3&    &    &     & 44  & 48\\
 \Delta^4&    &    &     &     & \mathbf{4}
\end{array}
\]
A function agreeing on five consecutive integers with a polynomial of degree at
most $3$ has vanishing fourth forward difference there. Here it is $4\ne0$, and
the window $n=5,\dots,9$ lies entirely inside the asserted range $n\ge5$.
\end{proof}

\begin{remark}
The five integers of \eqref{eq:witness} are not a printing accident: they are
forced by \eqref{eq:P} from other cells of the same table,
\begin{align*}
 v_5^3&=8+4+0+0-0=12,\\
 v_6^5&=36+36+12+8-4=88,\\
 v_7^7&=80+128+88+36-36=296,\\
 v_8^9&=140+292+296+80-128=680,\\
 v_9^{11}&=216+544+680+140-292=1288 .
\end{align*}
The cells consumed here are recorded in Table~\ref{tab:cells}.
\end{remark}

\begin{table}[ht]
\caption{Cells of the degree-distribution table of~\cite{CHOPTY} quoted in
this note. A dash is a printed blank, read as $0$. Rows are indexed by the
degree $d$ and columns by $n=2,\dots,9$.}
\label{tab:cells}
\[
\begin{array}{r|rrrrrrrr}
 d\backslash n & 2 & 3 & 4 & 5 & 6 & 7 & 8 & 9\\\hline
 3  & 0 & 4 & 8 & 12 & 16 & 20 & 24 & 28\\
 5  & - & 0 & 8 & 36 & 88 & 168 & 280 & 428\\
 7  & - & - & 0 & 12 & 80 & 296 & 792 & 1744\\
 9  & - & - & - & 0 & 16 & 140 & 680 & 2396\\
 11 & - & - & - & - & 0 & 20 & 216 & 1288
\end{array}
\]
\medskip
Additionally: the whole $n=3$ column, $v_3^2=4$, $v_3^3=4$, $v_3^4=8$,
$v_3^5=0$, $v_3^6=2$; the $a=6$ reverse diagonal
$v_5^4,v_6^6,v_7^8,v_8^{10},v_9^{12}=36,128,292,544,900$; and $v_4^2=4$ from
the convention at line 1204.
\end{table}

\section{The infinite family}

\subsection*{The reverse-diagonal recurrence}
Substituting $d=2n-a$ into \eqref{eq:P} and using $f_a(n)=v_n^{2n-a}$,
\begin{equation}\label{eq:R}
 f_a(n)=f_a(n-1)+f_{a-1}(n-1)-f_{a-1}(n-2)+f_{a-2}(n-1)+f_{a-2}(n-2)
\end{equation}
for every $n\ge3$.
Put $F_a(x):=\sum_{n\ge1}f_a(n)x^n$. Multiplying \eqref{eq:R} by $x^n$ and
summing over $n\ge3$ leaves the boundary polynomial
$f_a(1)x+\bigl[f_a(2)-f_a(1)-f_{a-1}(1)-f_{a-2}(1)\bigr]x^2$. Now
$f_a(1)=v_1^{2-a}=0$ for $a\ge1$ and $f_a(2)=v_2^{4-a}=0$ for $a\ge3$, so
\emph{for every $a\ge3$ the boundary vanishes identically} and
\begin{equation}\label{eq:F}
 (1-x)F_a \;=\; x(1-x)F_{a-1}\;+\;x(1+x)F_{a-2}\qquad(a\ge3),
\end{equation}
with the two seeds, obtained from the $a=2$ instance of \eqref{eq:R} whose
boundary term $2x^2$ is not zero,
\[
 F_0=\frac{2x}{1-x},\qquad F_1=0,\qquad
 F_2=\frac{4x^2}{(1-x)^2},\qquad F_3=\frac{4x^3}{(1-x)^2}.
\]
Equivalently $f_2(n)=4(n-1)$ and $f_3(n)=4(n-2)$ exactly, which reproduces the
$a=2$ and $a=3$ rows of the source's reverse-polynomial data table.

\subsection*{The normalised numerator}
Write $k_a:=\lfloor a/2\rfloor$ and
\[
 G_a(x)\;:=\;(1-x)^{k_a+1}F_a(x).
\]
Since $k_a-k_{a-1}=1$ and $k_a-k_{a-2}=1$ for even $a$, while $k_a-k_{a-1}=0$
and $k_a-k_{a-2}=1$ for odd $a$, dividing \eqref{eq:F} by
$(1-x)^{k_a+1}$ leaves no remainder:
\begin{equation}\label{eq:G}
\begin{aligned}
 G_a&=x(1-x)G_{a-1}+x(1+x)G_{a-2} && (a\ \text{even}),\\
 G_a&=x\,G_{a-1}+x(1+x)G_{a-2} && (a\ \text{odd}),
\end{aligned}
\end{equation}
with $G_2=4x^2$ and $G_3=4x^3$. Hence every $G_a$, $a\ge2$, lies in
$\mathbb Z[x]$; explicitly
\begin{align*}
 G_4&=4x^3+8x^4-4x^5, & G_5&=8x^4+12x^5-4x^6,\\
 G_6&=4x^4+20x^5+8x^6-20x^7+4x^8, &&\\
 G_7&=12x^5+40x^6+16x^7-24x^8+4x^9, &&
\end{align*}
\[
 G_8=4x^5+36x^6+56x^7-36x^8-56x^9+32x^{10}-4x^{11}.
\]

\subsection*{The dictionary}
For any $F=\sum_{n}f(n)x^n$ with $f(n):=0$ for $n\le0$, and any $j\ge0$,
\begin{equation}\label{eq:dict}
 [x^n](1-x)^jF(x)=(\Delta^jf)(n-j),
\end{equation}
because $[x^n](1-x)^jF=\sum_i(-1)^i\binom ji f(n-i)$, and re-indexing $i=j-t$
turns this into $(-1)^j\sum_t(-1)^t\binom jt f(n-j+t)=(\Delta^jf)(n-j)$. With
$j=k_a+1$ this reads
\begin{equation}\label{eq:dictG}
 [x^n]G_a=\bigl(\Delta^{k_a+1}f_a\bigr)(n-k_a-1),
\end{equation}
i.e.\ $G_a$ \emph{is} the $(k_a+1)$-st difference table of $f_a$, shifted.

\subsection*{The degree and the leading coefficient}
\begin{proposition}\label{prop:deg}
For $m\ge1$, $\deg G_{2m}=3m-1$ and $\deg G_{2m+1}=3m$, and the leading
coefficient of $G_a$ is $L_a=4(-1)^{\lfloor a/2\rfloor+1}$ for every $a\ge2$.
Moreover $G_a(1)>0$ for every $a\ge2$.
\end{proposition}

\begin{proof}
Let $D_a=\deg G_a$. In \eqref{eq:G} the competing leading terms are
$-L_{a-1}x^{D_{a-1}+2}$ against $L_{a-2}x^{D_{a-2}+2}$ for even $a$, and
$L_{a-1}x^{D_{a-1}+1}$ against $L_{a-2}x^{D_{a-2}+2}$ for odd $a$. From
$D_2=2$, $D_3=3$ the claimed formula follows by induction, and in both parities
the first candidate \emph{strictly} dominates: for even $a$ one has
$D_{a-1}=D_{a-2}+1$, hence $D_{a-1}+2>D_{a-2}+2$; for odd $a$ one has
$D_{a-1}+1=D_{a-2}+3>D_{a-2}+2$. So no cancellation can occur, and
$L_a=-L_{a-1}$ for even $a$, $L_a=L_{a-1}$ for odd $a$. With $L_2=L_3=4$ this
gives $L_a=4(-1)^{\lfloor a/2\rfloor+1}$; in particular $|L_a|=4$ for all
$a\ge2$. Setting $x=1$ in \eqref{eq:G} kills the $x(1-x)$ term, so
$G_a(1)=2G_{a-2}(1)$ for even $a$ and $G_a(1)=G_{a-1}(1)+2G_{a-2}(1)$ for odd
$a$; with $G_2(1)=G_3(1)=4>0$, induction gives $G_a(1)>0$.
\end{proof}

\begin{corollary}\label{cor:threshold}
For every $a\ge2$: $f_a$ agrees with a polynomial of degree exactly
$\lfloor a/2\rfloor$ on $[a-1,\infty)$, and $a-1$ is the least such threshold.
The last nonvanishing value of $\Delta^{k_a+1}f_a$ is at base $n=a-2$, where it
equals $4(-1)^{\lfloor a/2\rfloor+1}$.
\end{corollary}

\begin{proof}
By \eqref{eq:dictG}, $\Delta^{k_a+1}f_a(n)=[x^{n+k_a+1}]G_a$, which vanishes
exactly when $n+k_a+1>D_a$, i.e.\ when $n>D_a-k_a-1$. Proposition~\ref{prop:deg}
gives $D_a-k_a=a-1$ in both parities ($3m-1-m=2m-1$ and $3m-m=2m$), so
$\Delta^{k_a+1}f_a$ vanishes on $[a-1,\infty)$ and equals $L_a$ at $n=a-2$.
Vanishing of the $(k_a+1)$-st difference on a half-line is agreement there with
a polynomial of degree $\le k_a$; the degree is exactly $k_a$ because
$\Delta^{k_a}f_a$ is eventually the constant $G_a(1)\ne0$, again
by \eqref{eq:dictG} with $j=k_a$ and Proposition~\ref{prop:deg}. The leading
coefficient of that polynomial is $G_a(1)/k_a!$.
\end{proof}

\begin{theorem}\label{thm:family}
The reverse-diagonal theorem of~\cite{CHOPTY} is false for every $a\ge7$. The
base indices $n$ inside the asserted range at which
$\Delta^{\lfloor a/2\rfloor+1}f_a(n)\ne0$ all lie in $\{s(a),\dots,a-2\}$, and
the largest of them is $n=a-2$, which carries the value
$4(-1)^{\lfloor a/2\rfloor+1}$. For $2\le a\le 6$, by contrast, every
nonvanishing value of $\Delta^{\lfloor a/2\rfloor+1}f_a$ lies strictly below
$s(a)$, so on the range $n\ge s(a)$ granted by the source's exclusion clause
those five indices satisfy the theorem. Taken literally, with first admissible
index $\lceil a/2\rceil+1$, the theorem fails at $a=6$ as well; see Section~1.
\end{theorem}

\begin{proof}
By Corollary~\ref{cor:threshold} the obstruction at $n=a-2$ is nonzero for
every $a\ge2$, and it sits inside the range $n\ge s(a)$ precisely when
\[
  a-2\ \ge\ \Big\lfloor\frac a2\Big\rfloor+2
  \iff \Big\lceil\frac a2\Big\rceil\ \ge\ 4
  \iff a\ge 7 .
\]
For those $a$ the window $n=a-2,\dots,a-2+k_a+1$ lies in the asserted range,
and $\Delta^{k_a+1}f_a(a-2)\ne0$ contradicts agreement there with a polynomial
of degree $\le\lfloor a/2\rfloor$. Conversely for $2\le a\le6$ the same
inequality shows $a-2<s(a)$, so every nonvanishing value of
$\Delta^{k_a+1}f_a$ lies strictly below the first admissible $n$ and the
statement holds on the whole of $n\ge s(a)$.
\end{proof}

\begin{remark}
This explains the source's own confirming window rather than merely
contradicting it. The caption of its reverse-polynomial data table (line 1250)
says: ``The code computed all sequences up to $a=16$, but only confirmed
polynomial patterns up to $a=6$.'' The obstruction exists for every $a\ge2$,
but $a-2<\lfloor a/2\rfloor+2$ exactly when $a\le6$, so for those $a$ it lies
below $s(a)$ and is invisible to any check confined to the range. The two
cut-offs are the same inequality $\lceil a/2\rceil\ge4$.
\end{remark}

\section{What survives, and the repaired hypothesis}

Proposition~\ref{prop:deg} and Corollary~\ref{cor:threshold} prove the half of
the reverse-diagonal theorem that the source only asserted:

\begin{corollary}
For every $a\ge2$, $f_a$ agrees on a half-line with a polynomial of degree
exactly $\lfloor a/2\rfloor$ and leading coefficient
$G_a(1)/\lfloor a/2\rfloor!$. The same holds for $a=0$ by inspection, since
$f_0(n)=v_n^{2n}=2$.
\end{corollary}

Consequently, replacing the range clause of that theorem by
\begin{equation}\label{eq:repair}
 n\ \ge\ \max\Bigl(\Bigl\lfloor \frac a2\Bigr\rfloor+2,\ a-1\Bigr)
\end{equation}
makes it true for every $a\ge2$, and for $a=0$; the degenerate index $a=1$ is
outside anything claimed here. The repair also supersedes the even-$a$
boundary question, since $\lfloor a/2\rfloor+2$ excludes the degree-$2$ cell in
both parities. At $a=7$ the repaired bound is $6$; the two terms record the
admissible start \eqref{eq:s} and the agreement threshold of
Corollary~\ref{cor:threshold} respectively. In particular
the weaker repair
$n\ge\lfloor a/2\rfloor+2$ is \emph{not} correct: it is $5$ at $a=7$, which is
the very threshold Theorem~\ref{thm:a7} refutes.

\medskip
\noindent\textbf{Status.}
What is refuted is the \emph{range clause} of the reverse-diagonal theorem, as
a statement quantified over all $a\in\N_0$ and all $n\ge\lceil a/2\rceil+1$.
Polynomiality and the degree $\lfloor a/2\rfloor$ are \emph{not} refuted: they
are true for every $a\ge2$ and are proved above, where the source asserts them
without proof. Nothing here touches the source's proved sibling theorem at lines
1208--1226 (fixed degree $d$), which our program corroborates at $m=2$ for
$d=2,\dots,40$ on a sampled window $n=\lceil d/2\rceil+1,\dots,85$ of its stated
range, finding no obstruction there. The defect is specific to the
reverse diagonal, whose range clause is the sibling's word for word with $d$
replaced by $a$.

\medskip
\noindent\textbf{What is not settled.}
\begin{itemize}
\item The source's recurrence proposition is \emph{not} re-proved here. It is
proved in~\cite{CHOPTY} at lines 1191--1202 and taken as given; what we add is
the $w$-elimination \eqref{eq:P} and the boundary conventions, which are pinned
against the source's own printed integers (Table~\ref{tab:cells}) rather than
re-derived from the degree-extension-forest construction. If that proposition
failed at a boundary not exercised by those cells, \eqref{eq:R} would need
redoing.
\item No stage of this work brute-forced mountain--valley assignments of
$M_{2,n}$; the flat-foldability model is trusted, legitimately, because the
source proves the recurrence it comes from. Theorem~\ref{thm:a7} does not
depend on the model at all, since its five values are printed cells. The family
$a\ge8$, and every value at $n\ge10$, does.
\item Nothing is claimed about the theorem at the degenerate index $a=1$. The
seed $F_1=0$ above records that the reverse diagonal there is the zero sequence
(the census reproduces $f_1(n)=0$ for $n\le39$), so whether the asserted
``degree exactly $\lfloor1/2\rfloor=0$'' is to count as failing at $a=1$ turns
on the convention assigning a degree to the zero polynomial and on nothing else.
We draw no conclusion from it, and the refutation does not use it.
\item Beyond $a=80$ and $n=130$ nothing is checked by computation.
Proposition~\ref{prop:deg}, Corollary~\ref{cor:threshold} and
Theorem~\ref{thm:family} are proved for all $a$; the accompanying census only
corroborates them inside those bounds. Which indices of
$\{s(a),\dots,a-2\}$ other than $n=a-2$ actually fail is not asserted here.
\end{itemize}

\medskip
\noindent\textbf{Related work, and a prior threshold of the same shape.}
Gupta~\cite{GuptaOne} enumerates degree-$4$ and degree-$5$ vertices of
$\OFG(M_{m,n})$ by height functions, together with diameters; it says nothing
about a reverse diagonal (its ``antidiagonal'' is the grid antidiagonal
$\{i+j=s\}$ of a height function). Gupta~\cite{GuptaTwo} reached, in the
\emph{fixed-$d$ bivariate} slicing of the same family, a validity region
$m,n\ge\max(d-1,2)$ with the count departing from the polynomial below it, and
an onset magnitude $4$: ``one step below, this departure has leading
coefficient $-4$ times a Baxter number'' through $d=11$. So neither the
threshold shape ``parameter minus one'' nor the magnitude $4$ is presented here
as unprecedented. Neither work settles the reverse-diagonal slicing $d=2n-a$:
\cite{GuptaTwo} contains no such slicing, at $m=2$ its region
$m,n\ge\max(d-1,2)$ forces $d\le3$, and it treats the $m=2$ results
of~\cite{CHOPTY} as correct and settled.

\section{Verification}

Theorem~\ref{thm:a7} needs no computer: five integers printed
in~\cite{CHOPTY} and ten subtractions.
Theorem~\ref{thm:family}, Proposition~\ref{prop:deg},
Corollary~\ref{cor:threshold} and the repair \eqref{eq:repair} are hand algebra,
printed above in full with $G_2,\dots,G_8$ written out.

The accompanying program \texttt{verify.py} (Python~3.9+, standard library only,
exact integer arithmetic throughout) re-derives the \emph{finite} assertions of
this note from the objects printed above: the seeds $v_1^2=2$, $v_2^2=4$,
$v_2^3=0$, $v_2^4=2$; the recurrence \eqref{eq:P}; the cells of
Table~\ref{tab:cells}; the witness \eqref{eq:witness}; the polynomials
$G_2,\dots,G_8$. It calibrates the recurrence against the source's printed
integers before using it, verifies the dictionary \eqref{eq:dictG} against the
recurrence, checks the degree, leading-coefficient and positivity statements of
Proposition~\ref{prop:deg}, the threshold $a-1$ and the repair
\eqref{eq:repair} for $a\le80$ --- both from $G_a$ and directly from difference
tables --- and censuses the source's proved sibling theorem as a null. It does
\emph{not} verify the universal arguments: the proofs of
Proposition~\ref{prop:deg}, Corollary~\ref{cor:threshold} and
Theorem~\ref{thm:family} are hand proofs and are only corroborated there; the
source's recurrence proposition is not re-proved; no mountain--valley assignment
is enumerated anywhere; and the line numbers, byte count and md5 of the source
file are not fetched. Its transcript, with its own closing statement of scope,
is \texttt{verify.output.txt}.

\begin{thebibliography}{9}

\bibitem{CHOPTY}
L.~Christensen, T.~C. Hull, E.~O'Neil, V.~Pappano, N.~Ter-Saakov, and K.~Yang,
\emph{The origami flip graph of the $2\times n$ Miura-ori},
\href{https://arxiv.org/abs/2506.19700v2}{arXiv:2506.19700v2}, 2026.
Quotations and line numbers refer to the \LaTeX{} source
\texttt{Miuradraft-v2.tex} of the v2 e-print ($148\,285$ bytes, md5
\texttt{adca23f5c2eade3979cfc77682341c6d}); we located no later e-print version
and no journal reference for it.

\bibitem{GuptaOne}
C.~Gupta,
\emph{Height functions on the $m\times n$ Miura-ori flip graph: degree sequence
and diameter},
\href{https://arxiv.org/abs/2606.22614}{arXiv:2606.22614v2}, 2026.

\bibitem{GuptaTwo}
C.~Gupta,
\emph{One construction for the Miura-ori flip-graph degree sequence},
\href{https://arxiv.org/abs/2607.05567}{arXiv:2607.05567v3}, 26 August 2026.

\end{thebibliography}

\end{document}
